\subsection{The bound the two tests are read against}

The surface analysis \cite{MMC2026b} gives, for the glued solve,
\begin{equation}
  \bigl\| Q_h u_h - u \bigr\|_\infty
  \;\le\; C\Bigl\{ h^{q} + h^{p-1} + h\,(D + J_2) + h^2
          + \varepsilon_R + h\,\rho_{f,h} \Bigr\} ,
  \label{eq:bound}
\end{equation}
conditional on a discrete stability assumption. With $q=2$ and $p=3$,
and with a source-consistent forcing convention, this reads
\begin{equation}
  \bigl\| Q_h u_h - u \bigr\|_\infty
  \;\le\; C\bigl\{ h^2 + h\,(D + J_2) + \varepsilon_R \bigr\} .
  \label{eq:bound2}
\end{equation}
Three constants in \eqref{eq:bound2} decide what the tables can show.

\paragraph{$D$, the extrinsic junction defect.} $D$ is the supremum over
the junction curve of $|b_2^\perp - b_1^\perp|$, the mismatch of only the
\emph{mixed} and \emph{transverse} parts $(b_{sr}, b_{rr})$ of the second
fundamental form in the matched edge frame. On the flipped ellipsoid the
meridian and the parallel are the principal directions of both branches,
so $b_{sr}=0$ on both and
\begin{equation}
  D = 2\kappa_m .
\end{equation}
The along-edge curvature $\kappa_p$ drops out of $D$, because the
transmission conditions already match $U_{ss}$.

\paragraph{$\dk$, the quantity the tables report.} The ellipsoid script
reports instead
\begin{equation}
  \dk = \bigl\| \mathrm{II}_1 - R\,\mathrm{II}_2 R^{\!\top} \bigr\|_F ,
\end{equation}
the full matched-frame mismatch, which \cite{MMC2026b} calls a valid but
less sharp replacement for $D$. On this family $\dk$ runs 18 to 34 per
cent above $D$, and it is monotone in the same direction, so ranking the
cuts by either gives the same order.

The script builds $\dk$ rather than assuming it. Both branches are pieces
of the same ellipsoid, met at the same $\varphi_c$, so they share
\begin{equation}
  \kappa_m = \frac{ac}{G^{3/2}}, \qquad
  \kappa_p = \frac{c}{a\sqrt{G}}, \qquad
  G(\varphi) = a^2\cos^2\!\varphi + c^2\sin^2\!\varphi ,
\end{equation}
and the mismatch collapses to $|n_1 - Rn_2|\sqrt{\kappa_m^2+\kappa_p^2}$
with $n$ the inward unit normals. On the \emph{unflipped} cut, $R$ is the
identity and $n_2=n_1$, so $\dk$ vanishes identically and the glued solve
should recover the second-order rate of the uncut one. On the
\emph{flipped} cut, $R$ carries $n_2$ to $-n_1$, so $|n_1-Rn_2|=2$ and
\begin{equation}
  \dk = 2\sqrt{\kappa_m^2 + \kappa_p^2}
  \qquad\text{(flipped)}
\end{equation}
at every cut height. It varies only through where on the meridian the cut
falls.

\paragraph{$J_2$, the intrinsic junction defect.} $J_2$ is
$[U_{rr}] = (k_2-k_1)U_r - [f]$, with $k_i$ the geodesic curvature of the
junction curve inside branch $i$. On the flipped ellipsoid it vanishes,
and \emph{both} of its terms vanish separately. The two branches are
mirror images across the cut plane, so $k_2=k_1$. And $u$ and $f$ are the
same functions of $(\varphi,\vartheta)$ on both branches, so $[f]=0$.

That is a different route to $J_2=0$ than the two-sphere example of
\cite{MMC2026b}, where the two terms are individually nonzero and cancel.
On the flipped ellipsoid the whole junction defect is therefore
\emph{extrinsic}, and the bound reads
\begin{equation}
  \bigl\| Q_h u_h - u \bigr\|_\infty
  \;\le\; C\bigl\{ h^2 + h\,D + \varepsilon_R \bigr\} .
  \label{eq:bound3}
\end{equation}

\paragraph{$\varepsilon_R$, the rotation error.} This is the term the
estimated rotations pay and the exact rotation does not. It is the only
term that separates the three variants of the ellipsoid test, and the
only term that separates \texttt{exact} from \texttt{dk2} on the tube.

\paragraph{$\rho_{f,h}$, the forcing residual.} This is why both scripts
take $f$ at the \emph{rotated target} on a routed row. Evaluating the
forcing on the opposite side without correcting a trace jump can leave
$\rho_{f,h}=O(1)$ instead of $O(h)$. That would promote $h\rho_{f,h}$
from $O(h^2)$ to $O(h)$ and put it on a level with the junction term
itself. It is safe in both tests only because $[f]=0$ by construction.
